\section{How the three empirical solvers fit}
\label{sec:empirical-map}

The framework separates four roles that the measured notes currently mix.

\begin{center}
\small
\begin{tabularx}{\textwidth}{@{}p{0.20\textwidth}p{0.21\textwidth}p{0.23\textwidth}X@{}}
\toprule
Method & Propagation & Settlement & Revisit control \\
\midrule
Two-rung SOR & one ranked $\omega_\star$ phase above $2.5g$ & repeated
$\omega=1$ coordinate pushes at $g$ & two declared regions; no feedback \\
Alternating ladder & several over-relaxed bands & interleaved unit-relaxation
bands & fixed threshold schedule \\
Adaptive frontier ladder & over-relaxed alternating bands with stacked
guards & interleaved unit-relaxation bands & measured work divided by touched
support, transferred between related instances \\
Ideal theoretical successor & causal directed messages & one exact block or
certified inexact block solve & control $\mathfrak R_{\rm set}$ or the
separator/signature response complexity \\
\bottomrule
\end{tabularx}
\end{center}

The adaptive statistic is not literally
$\mathfrak R_{\rm set}$: it counts executed row operations and uses the final
touched support in the denominator.  Under the bounded-scan hypotheses of
\cref{thm:propagate-settle-work}, however, it is an observable constant-factor
proxy for the same phenomenon.  This suggests a principled role for feedback:
adapt the propagation strength or trigger settlement when the observed
operation revisit factor begins to exceed the theoretical budget.

The empirical band factor $2.5$ is a finite-parameter discovery rule, not the
asymptotic spacing in the accelerated theorem.  Uniform no-reactivation for
optimal SOR as $\alpha\downarrow0$ requires an alpha-scaled factor
$1+\Theta(\sqrt\alpha)$, or directed state that makes reactivation irrelevant.
